%% ---------------------------------------------------------------------
%% 5. TECHNICAL DESCRIPTION — Required · 1–5 pages
%% The engineering / scientific content the student needs to start working.
%% For large projects split it into several chapter files
%% (05a-hardware.tex, 05b-communication.tex, 05c-data.tex …) and \input each in main.tex.
%% Typical content:
%%  - components and their responsibilities; interfaces and data formats;
%%  - key parameters, register maps, protocols, data models (tables);
%%  - reference designs / existing code the student starts from;
%%  - known difficulties and hints.
%% Describe WHAT is required and the constraints; leave design decisions the
%% student is expected to make open, and say so ("the student chooses …").
%% ---------------------------------------------------------------------
\chapter{Technical Description}
\label{ch:technical}

\section{Components}
\missing{Describe each component and its responsibility.}

\section{Interfaces and data}
\missing{Describe interfaces, protocols and data formats. Use a table for parameters / registers.}

\begin{table}[htbp]
  \centering
  \caption{Example parameter table (replace).}
  \label{tab:params}
  \begin{tabular}{lll}
    \toprule
    \textbf{Parameter} & \textbf{Range} & \textbf{Note} \\
    \midrule
    \missing{name} & \missing{range} & \missing{note} \\
    \bottomrule
  \end{tabular}
\end{table}

\section{Starting point}
\missing{Existing code, reference designs, lab materials the student receives on day one.}
